\documentclass[10pt,a4paper]{amsart}
\usepackage[margin=19mm]{geometry}
\usepackage{amsmath,amssymb,mathtools}
\usepackage[scr=rsfs,cal=euler]{mathalfa}
\usepackage{xfrac}
\usepackage[unicode,hidelinks,bookmarks=false]{hyperref}
\newcommand{\ie}{i.e.\ }
\newcommand{\cf}{cf.\ }
\newcommand{\resp}{respectively\ }
\newcommand{\Subsection}{\subsection}
\providecommand{\nobreakdash}{\nobreak\hbox{-}}
\setlength{\emergencystretch}{2em}
\multlinegap0pt
\usepackage{amssymb}
\usepackage{enumerate}
\usepackage{bm}
\usepackage{mathtools}
\usepackage{float}
\usepackage{algorithm}\usepackage{algpseudocode}
\newtheorem{theorem}{Theorem}
\newcommand {\C}     {\mathrm{C}}
\newcommand{\End}{\mathrm{End}}
\newcommand{\card}{\mathrm{card}}
\newcommand{\mmrank}{\mathrm{mmrank}}
\newcommand{\rank}{\mathrm{rank}}
\title{Bounds on the realizations of zero-nonzero patterns and sign conditions of polynomials restricted to varieties and applications}
\author{Saugata Basu, Laxmi Parida}
\date{Source: arXiv:2411.11729v5 (CC BY 4.0)}
\begin{document}
\maketitle

\noindent\emph{Source and licence.} Bounds on the realizations of zero-nonzero patterns and sign conditions of polynomials restricted to varieties and applications. Version arXiv:2411.11729v5 (CC BY 4.0), distributed by arXiv under the Creative Commons Attribution 4.0 International licence, http://creativecommons.org/licenses/by/4.0/, as recorded in the arXiv licence metadata for this exact version and shown on the abstract page https://arxiv.org/abs/2411.11729v5. Full licence notice: \texttt{licenses/arxiv-2411.11729-CC-BY-4.0.txt}.\par\smallskip

\noindent\textit{Editorial scope.} Counted unit: the article's theorem thm:quantum2 (source0.tex lines 3806--3813). The article's complete original proof of it is delivered with it (source0.tex lines 4232--4386, one \texttt{proof} environment of 155 lines, which the article places away from the statement). No mathematics is written, repaired or re-derived: the statement and the proof are the article's own lines, in the article's own notation, and no retained line is substituted or re-flowed.  No further source-local context is needed: every cross-reference of every retained line resolves inside the two delivered spans.  Nothing was omitted from any retained span, so the package carries no omission record.  No retained line contains a citation, so the wrapper carries no bibliography and no bibliography entry.  No retained line contains a figure environment, an image-inclusion command or drawing code, so no image is delivered and the package declares no asset dependency.  Editorial formatting only, in addition: an AMS \texttt{amsart} preamble that restates the article's own declarations and macros used by the retained lines, namely the environments \texttt{theorem} on separate counters, together with the standard packages those lines need.  The retained environments print in this document as Theorem 1 in order of appearance, whereas the article prints the same items with the numbering of its own full document.  The complete native source of the article as distributed in the arXiv e-print for this exact version is bundled unchanged as \texttt{sources/168694/source0.tex}.
\begin{theorem}
\label{thm:quantum2}
\label{cor:15}
Let $n>0, t \geq 0$, and  $A=\End(V^{\otimes(n+t+1)})$, viewed as a real algebra, and let $\Delta\subset A_\C$ be a real algebraic set whose real points are equidimensional on each irreducible component, with $\dim\Delta=p$ and $\deg\Delta=D$. Then there is an absolute constant $C>0$ and a Boolean function $f:\{0,1\}^n\to\{0,1\}$ whose quantum complexity with $t+1$ ancillary qubits relative to $\Delta$ is at least
\[
C\cdot\frac{2^n}{(p+1)(n+\log D)}.
\]
\end{theorem}

\begin{proof}[Proof of Theorem~\ref{thm:quantum2}]
Put

$$
A=\End(\mathcal H_{n,t}),
$$

viewed as a real algebra. For each \(x\in\{0,1\}^{n}\), decompose

$$
E\lvert x\,0\,0^{t}\rangle
=
u_{x,0}(E)\oplus u_{x,1}(E)
$$

according to the value \(0\) or \(1\) of the output qubit. Thus

$$
u_{x,b}(E)\in
V^{\otimes n}\otimes \lvert b\rangle\otimes V^{\otimes t},
\qquad b\in\{0,1\},
$$

and

$$
Q_x(E)
=
\left\|u_{x,1}(E)\right\|^{2}
-
\left\|u_{x,0}(E)\right\|^{2}.
$$

We first observe that each \(\mathcal E^{+}_{f,t}\) is semi-algebraically
connected. Indeed, for a fixed \(x\), the condition \(Q_x(E)>0\) has the
form

$$
\|u_1\|^{2}>\|u_0\|^{2}.
$$

The homotopy

$$
(u_0,u_1)\longmapsto ((1-s)u_0,u_1),
\qquad 0\leq s\leq1,
$$

preserves this inequality and retracts the set onto

$$
\{0\}\times
\left(
V^{\otimes n}\otimes\lvert1\rangle\otimes V^{\otimes t}
\setminus\{0\}
\right),
$$

which is semi-algebraically connected. The same argument, with the two
factors interchanged, applies to \(Q_x(E)<0\).

The conditions corresponding to distinct \(x\)'s involve distinct columns
of the matrix of \(E\), while all other columns are unrestricted.
Consequently, \(\mathcal E^{+}_{f,t}\) is a product of semi-algebraically
connected sets and is therefore semi-algebraically connected.

Now let

$$
W^{+}
=
\bigcup_{f:\{0,1\}^{n}\to\{0,1\}}
\mathcal E^{+}_{f,t}.
$$

Equivalently,

$$
W^{+}
=
\left\{
E\in A
\;\middle|\;
Q_x(E)\neq0
\text{ for every }x\in\{0,1\}^{n}
\right\}.
$$

If \(f\neq g\), then for some \(x\) the sets
\(\mathcal E^{+}_{f,t}\) and \(\mathcal E^{+}_{g,t}\) impose opposite
strict signs on \(Q_x\), and hence are disjoint. Since each
\(\mathcal E^{+}_{f,t}\) is nonempty and semi-algebraically connected,
these sets are precisely the semi-algebraically connected components of
\(W^{+}\). Therefore

$$
b_0(W^{+})=2^{2^{n}}.
$$

Let

$$
\mathcal P=\{Q_x\mid x\in\{0,1\}^{n}\}.
$$

Then \(W^{+}\) is a \(\mathcal P\)-semi-algebraic set,

$$
\card(\mathcal P)=2^{n},
\qquad
\deg Q_x\leq2.
$$

Applying Theorem~13 gives

$$
\mmrank_{A,\Delta}(W^{+})
\geq
C\,
\frac{2^{n}}
     {(p+1)(n+\log D)}
$$
for an absolute constant \(C>0\).
Since the semi-algebraically connected components of \(W^{+}\) are exactly
the sets \(\mathcal E^{+}_{f,t}\), there exists a Boolean function \(f\)
such that

$$
\min_{E\in\mathcal E^{+}_{f,t}}
\rank_{A,\Delta}(E)
\geq
C\,
\frac{2^{n}}
     {(p+1)(n+\log D)}.
$$

Finally,

$$
\mathcal U^{+}_{f,t}\subset\mathcal E^{+}_{f,t},
$$

and hence

$$
\min_{U\in\mathcal U^{+}_{f,t}}
\rank_{A,\Delta}(U)
\geq
\min_{E\in\mathcal E^{+}_{f,t}}
\rank_{A,\Delta}(E).
$$

The left-hand side is the quantum complexity of \(f\) with \(t+1\)
ancillary qubits relative to \(\Delta\), which proves the theorem.
\end{proof}

\end{document}
